\chapter*{Abstract}
\addcontentsline{toc}{chapter}{Abstract}
\begin{center}\large\bfseries\thesistitle\end{center}

Local surface rework requires coordinating tool movement, contact force, and time within a selected region. This thesis studies a setting in which a person specifies the region and demonstrations supply its temporal processing strategy. Rules and demonstration replay are valid alternatives; the evaluation examines the execution characteristics of observation-conditioned skill generation and its reuse at registered planar placements.

The framework connects temporally paired tool pose, three-axis interaction force, and RGB observations to conditional Flow Matching. A joint motion--force action sequence represents position, orientation, and teaching-derived target force; measured force and its history separately condition the policy. Position is learned relative to the stain center. A direction adapter aligns observation XY and the policy image, then inversely restores generated XY. A shared executor conditions the references before existing lower-level force control regulates contact.

Three evaluations address strategy, force information, and spatial reuse. E1's latest summary separates five additional rule runs (\(R^{*}\), 19 October) from original T/C runs (15 October), preserving original R. Each completes the contact--loaded-pass--exit procedure in its observed five runs. Mean times are \EoneRstarTime, \EoneTTime, and \EoneCTime~s for \(R^{*}\)/T/C; T has smaller time dispersion and excess-force integral than C. The additional rule retiming reduces mean time by \RstarTimePercent\% relative to original R, without matching measured mean speed. E2 distinguishes pose-only A from pose--force B/C with masked or measured-force conditioning. A completes four of five attempts, retaining A4's interrupted release transition and null completion time. B/C complete five each; C's mean operator-based interval is \BCSeconds~s (\BCPercent\%) shorter than B's. These descriptive intervals differ from E1's automatic procedure times.

E3 reuses the same C checkpoint and normalizer without new teaching or retraining. Saved-image analysis measures board-marker area reduction using a common gray threshold and saved visible-boundary annotations. N/P1/D1/D2 contain \EthreeMainN{} selected runs and \EthreeImageN{} image-evaluable pairs. Removal means \(\pm\) sample SD are \EthreeNRemoval\(\pm\)\EthreeNSD, \EthreePoneRemoval\(\pm\)\EthreePoneSD, \EthreeDoneRemoval\(\pm\)\EthreeDoneSD, and \EthreeDtwoRemoval\(\pm\)\EthreeDtwoSD\%, with image \(n=5,5,4,5\). D1-4 remains selected with final tool occlusion. P2's five executed, occluded runs remain outside the main results in the 25-run archive. The findings support an executable human-guided skill, observed force-configuration differences, and reuse with registration and direction adaptation in the tested placements. Rules and replay remain effective fixed-task alternatives. E3 ink-area reduction is distinct from E1/E2 completion and does not establish independently measured metal surface-quality or micro-scratch restoration.

\vspace{0.4cm}
\noindent\textbf{Keywords:} Local Robotic Polishing, Human Task Specification, Motion--Force Skill, Imitation Learning, Flow Matching, Spatial Skill Transfer, Image-Based Stain Removal
